\documentclass[12pt]{article}
\usepackage{fontspec}
\usepackage{graphicx}
\usepackage{scrextend}
\usepackage{listings}
\usepackage{tikz}
\usepackage{amssymb}
\usepackage{color}
\usepackage{soul}
\usepackage{mathrsfs}
\usepackage{fancyhdr}
\usepackage{pgfornament}
\usepackage[many]{tcolorbox}
\usepackage{collectbox}
\usepackage{mdframed}
\usepackage{xcolor}
\usepackage{mathtools}
\usepackage[hidelinks]{hyperref}
\usepackage[final]{pdfpages}
\usepackage[left=0.5in,right=0.5in,top=0.5in,bottom=1in,footskip=15mm]{geometry}

\usepackage{xcolor}
\linespread{1.2}

\begin{document}
\pagenumbering{gobble}

\begin{center}
\begin{Large}
\textbf{Information about Imaginary Numbers, version 0} \\

\vspace{5mm}
\end{Large}
\end{center}

\vspace{10mm}

\begin{addmargin}[3em]{4em}
\textbf{Statement:} \hspace{2mm} The imaginary number $i$ is such that $i\neq0, i \not< 0$, and $i\not>0$.

\vspace{5mm}

\noindent\textbf{Proof:} Clearly $i\neq 0$ because $i\overset{\text{def}}{=}\sqrt{-1}$, so suppose $i<0$. If we multiply this inequality by $i$ on both sides, we obtain $$i\cdot i = i^2 = -1 > i \cdot 0 = 0,$$ or simply $-1 > 0$ (note that the direction of the inequality symbol flipped because we supposed $i<0$). This is a contradiction. Now suppose $i>0$. Multiplying this inequality by $i$ on both sides yields $$i\cdot i = i^2 = -1 > i\cdot 0 = 0,$$ i.e., $-1>0$, the same contradiction as before. Thus, we conclude that $i\neq 0$, $i\not<0$, and $i\not>0.$

\hfill$\blacksquare$

\vspace{80mm}

\noindent\textbf{Euler's Identity (``Euler'' is pronounced like ``Oiler''):} It can be shown using calculus that $$\sin(x) = x - \frac{x^3}{3!} + \frac{x^5}{5!} - \frac{x^7}{7!} + \ldots$$ $$\cos(x) = 1 - \frac{x^2}{2!} + \frac{x^4}{4!} - \frac{x^6}{6!} + \ldots$$ $$e^x = 1 + x + \frac{x^2}{2!} + \frac{x^3}{3!} + \frac{x^4}{4!} + \frac{x^5}{5!} + \ldots$$ 

\vspace{2mm}

\noindent These are called Maclaurin Series. Let's try adding $\cos(x)$ and $\sin(x)$ together to see what happens. We get

\begin{align*}
\cos(x) + \sin(x) &= \left(1 - \frac{x^2}{2!} + \frac{x^4}{4!} - \frac{x^6}{6!} + \ldots\right) + \left( x - \frac{x^3}{3!} + \frac{x^5}{5!} - \frac{x^7}{7!} + \ldots\right) \\
&= 1 + x - \frac{x^2}{2!} - \frac{x^3}{3!} + \frac{x^4}{4!} + \frac{x^5}{5!} - \frac{x^6}{6!} - \frac{x^7}{7!} + \ldots
\end{align*}

\noindent Notice that this looks almost like the Maclaurin Series for $e^x$, but some of the signs are off. Can we fix that? Notice that all of the terms with powers that are multiples of $4$ have a positive sign in front of them. Furthermore, all the terms with even powers that are not multiples of $4$ have negative signs in front of them. This is consistent with how $i$ behaves; for example, $i^2 = -1$, $i^4=1$, $i^6=-1$, $i^8=1$ and so on. Therefore, we can write

\begin{align*}
\cos(x) + \sin(x) &= \left(1 - \frac{x^2}{2!} + \frac{x^4}{4!} - \frac{x^6}{6!} + \ldots\right) + \left( x - \frac{x^3}{3!} + \frac{x^5}{5!} - \frac{x^7}{7!} + \ldots\right) \\
&= 1 + x - \frac{x^2}{2!} - \frac{x^3}{3!} + \frac{x^4}{4!} + \frac{x^5}{5!} - \frac{x^6}{6!} - \frac{x^7}{7!} + \ldots \\
&= 1 + x + i^2 \frac{x^2}{2!} - \frac{x^3}{3!} + i^4 \frac{x^4}{4!} + \frac{x^5}{5!} + i^6 \frac{x^6}{6!} - \frac{x^7}{7!} + \ldots \\ 
&= 1 + x + \frac{(ix)^2}{2!} - \frac{x^3}{3!} + \frac{(ix)^4}{4!} + \frac{x^5}{5!} + \frac{(ix)^6}{6!} - \frac{x^7}{7!} + \ldots
\end{align*}

\noindent This makes all of the even terms positive. We still need to make all the terms with odd exponents positive to make this look like the Maclaurin Series for $e^x$. These terms have signs consistent with the odd powers of $i$. For example, $i^1=i$, $i^3=-i$, $i^5=i$, $i^7=-i$, and so on. While the signs are right, there are in fact no $i$'s in these terms. In order to get them in there, we have to actually multiply them by $i$. So let us look at $\cos(x) + i\sin(x).$ We have 

\begin{align*}
\cos(x) + i\sin(x) &= 1 + ix + \frac{(ix)^2}{2!} - i\frac{x^3}{3!} + \frac{(ix)^4}{4!} + i\frac{x^5}{5!} + \frac{(ix)^6}{6!} - i\frac{x^7}{7!} + \ldots \\
&= 1 + ix + \frac{(ix)^2}{2!} + \frac{(ix)^3}{3!} + \frac{(ix)^4}{4!} + \frac{(ix)^5}{5!} + \frac{(ix)^6}{6!} + \frac{(ix)^7}{7!} + \ldots \\
&= e^{ix}.
\end{align*}

\noindent Having established that $e^{ix} = \cos(x) + i\sin(x),$ we can now conclude that $$e^{i\pi} = \cos(\pi) + i \sin(\pi) = -1 + i \cdot 0 = -1.$$ In other words, $$e^{i\pi} + 1 = 0.$$

\end{addmargin}

\end{document}
